\documentclass[11pt,a4paper]{article}
\usepackage[T1]{fontenc}
\usepackage[utf8]{inputenc}
\usepackage[provide=*,lithuanian]{babel}
\usepackage{lmodern,amsmath,amssymb,enumitem}
\usepackage[margin=1.5cm]{geometry}
\setlength{\parindent}{0pt}
\setlength{\emergencystretch}{4em}
\setlength{\abovedisplayskip}{4pt}
\setlength{\belowdisplayskip}{4pt}
\setlist[enumerate]{label=\textbf{\arabic*.},leftmargin=0.7cm,itemsep=0.24cm,topsep=0.2cm,parsep=0pt}
\begin{document}
\begin{center}
{\Large\bfseries Vijetos teoremos taikymas}

\medskip
{\large Lygis: Boss}

\end{center}
\textbf{Užduotis.} Išspręskite uždavinius, remdamiesi Vijetos teorema
ir algebrinėmis tapatybėmis. Skaičiuotuvu nesinaudokite.
Uždaviniuose su parametru patikrinkite gautų reikšmių tinkamumą.
\medskip
\begin{enumerate}
\item \begin{minipage}[t]{\linewidth}Duoto daugianario \(P(x)=x^{2}-4x+2\) šaknys yra \(x_1,x_2\). Nerasdami šaknų, apskaičiuokite \mbox{\(\frac{1}{x_{2}^{3} + 1} + \frac{1}{x_{1}^{3} + 1}\)}.\end{minipage}
\item \begin{minipage}[t]{\linewidth}Raskite visas \(m\) reikšmes, su kuriomis lygties \(x^{2}-mx+m - 2=0\) realiosios šaknys \(x_1,x_2\) tenkina \mbox{\(x_1^3+x_2^3=18\)}.\end{minipage}
\item \begin{minipage}[t]{\linewidth}Lygties \(x^{2}-x-1=0\) šaknys yra \(x_1,x_2\). Nerasdami šaknų, apskaičiuokite \mbox{\(x_{1}^{5} + x_{2}^{5}\)}.\end{minipage}
\item \begin{minipage}[t]{\linewidth}Daugianario \(P(x)=x^{2}-mx+m - 1\) realiosios nenulinės šaknys yra \(x_1,x_2\). Raskite visas \(m\) reikšmes, su kuriomis \mbox{\(\frac{x_1}{x_2}+\frac{x_2}{x_1}=4\)}.\end{minipage}
\item \begin{minipage}[t]{\linewidth}Raskite visas \(m\) reikšmes, su kuriomis daugianariai \[P(x)=x^{2}-\left(m + 1\right)x+m - 2,\qquad Q(x)=x^{2}-3x+1\] turi bent vieną bendrą realiąją šaknį.\end{minipage}
\item \begin{minipage}[t]{\linewidth}Raskite visas \(m\) reikšmes, su kuriomis daugianario \(P(x)=x^{2}-mx+m + 6\) abi šaknys yra sveikieji skaičiai.\end{minipage}
\item \begin{minipage}[t]{\linewidth}Kvadratinio daugianario \(P(x)\) vyriausiasis koeficientas lygus \(2\), o realiosios šaknys tenkina: \[x_1^2+x_2^2=13,\quad x_1^3+x_2^3=35.\] Raskite visus galimus daugianarius \(P(x)\).\end{minipage}
\item \begin{minipage}[t]{\linewidth}Lygties \(x^{2}-mx+1=0\) realiosios šaknys yra \(x_1,x_2\). Raskite visas \(m\) reikšmes, su kuriomis \mbox{\(x_1^2,\quad x_2^2\)} yra lygties \(x^{2}-\left(m + 2\right)x+1=0\) šaknys.\end{minipage}
\item \begin{minipage}[t]{\linewidth}Duoto daugianario \(P(x)=x^{2}-4x+2\) šaknys yra \(x_1,x_2\). Nerasdami jų, sudarykite kvadratinę lygtį, kurios vyriausiasis koeficientas lygus \(1\), o šaknys \mbox{\(x_{1}^{2} + x_{2},\quad x_{1} + x_{2}^{2}\)}.\end{minipage}
\item \begin{minipage}[t]{\linewidth}Duoto daugianario \(P(x)=x^{2}-5x+3\) šaknys yra \(x_1<x_2\). Nerasdami šaknų, apskaičiuokite \mbox{\(\frac{x_{1}^{2}}{x_{2}} - \frac{x_{2}^{2}}{x_{1}}\)}.\end{minipage}
\item \begin{minipage}[t]{\linewidth}Lygties \(x^{2}-mx+1=0\) realiosios šaknys yra \(x_1,x_2\). Raskite visas \(m\) reikšmes, su kuriomis \mbox{\(x_1^2,\quad x_2^2\)} yra lygties \(x^{2}-\left(m + 4\right)x+1=0\) šaknys.\end{minipage}
\item \begin{minipage}[t]{\linewidth}Lygties \(x^{2}-3x+1=0\) šaknys yra \(x_1,x_2\). Nerasdami jų, sudarykite kvadratinę lygtį, kurios vyriausiasis koeficientas lygus \(1\), o šaknys \mbox{\(x_{1}^{2} - x_{2},\quad - x_{1} + x_{2}^{2}\)}.\end{minipage}
\item \begin{minipage}[t]{\linewidth}Lygties \(x^{2}-3x+1=0\) šaknys yra \(x_1<x_2\). Nerasdami šaknų, apskaičiuokite \mbox{\(x_{1}^{3} - x_{2}^{3}\)}.\end{minipage}
\item \begin{minipage}[t]{\linewidth}Kvadratinio daugianario \(P(x)\) vyriausiasis koeficientas lygus \(1\), o realiosios šaknys tenkina: \[x_1^2+x_2^2=7,\quad x_1^3+x_2^3=10.\] Raskite visus galimus daugianarius \(P(x)\).\end{minipage}
\item \begin{minipage}[t]{\linewidth}Raskite visas \(m\) reikšmes, su kuriomis daugianario \(P(x)=x^{2}-mx+2 m + 1\) abi šaknys yra teigiamieji sveikieji skaičiai.\end{minipage}
\item \begin{minipage}[t]{\linewidth}Raskite visas \(m\) reikšmes, su kuriomis lygties \(x^{2}-\left(m + 1\right)x+m - 1=0\) realiosios šaknys \(x_1,x_2\) tenkina \mbox{\(x_1^3+x_2^3=8\)}.\end{minipage}
\item \begin{minipage}[t]{\linewidth}Lygties \(x^{2}-3x+1=0\) šaknys yra \(x_1,x_2\). Nerasdami šaknų, apskaičiuokite \mbox{\(\frac{x_{1}}{x_{1}^{2} + 1} + \frac{x_{2}}{x_{2}^{2} + 1}\)}.\end{minipage}
\newpage
\item \begin{minipage}[t]{\linewidth}Raskite visas \(m\) reikšmes, su kuriomis daugianariai \[P(x)=x^{2}-\left(m + 2\right)x+m,\qquad Q(x)=x^{2}-2x-1\] turi bent vieną bendrą realiąją šaknį.\end{minipage}
\item \begin{minipage}[t]{\linewidth}Duoto daugianario \(P(x)=x^{2}+x-1\) šaknys yra \(x_1,x_2\). Nerasdami šaknų, apskaičiuokite \mbox{\(x_{1}^{6} + x_{2}^{6}\)}.\end{minipage}
\item \begin{minipage}[t]{\linewidth}Daugianario \(P(x)=x^{2}-\left(m + 1\right)x+1\) realiosios nenulinės šaknys yra \(x_1,x_2\). Raskite visas \(m\) reikšmes, su kuriomis \mbox{\(\frac{x_1}{x_2}+\frac{x_2}{x_1}=7\)}.\end{minipage}
\item \begin{minipage}[t]{\linewidth}Duoto daugianario \(P(x)=x^{2}-2x-1\) šaknys yra \(x_1,x_2\). Nerasdami šaknų, apskaičiuokite \mbox{\(\frac{1}{x_{2}^{2} + 2 x_{2} + 2} + \frac{1}{x_{1}^{2} + 2 x_{1} + 2}\)}.\end{minipage}
\item \begin{minipage}[t]{\linewidth}Lygties \(x^{2}-mx-1=0\) realiosios šaknys yra \(x_1,x_2\). Raskite visas \(m\) reikšmes, su kuriomis \mbox{\(x_1^2,\quad x_2^2\)} yra lygties \(x^{2}-\left(2 m + 2\right)x+1=0\) šaknys.\end{minipage}
\item \begin{minipage}[t]{\linewidth}Raskite visas \(m\) reikšmes, su kuriomis daugianario \(P(x)=x^{2}-mx+2 - m\) abi šaknys yra neigiamieji sveikieji skaičiai.\end{minipage}
\item \begin{minipage}[t]{\linewidth}Daugianario \(P(x)=x^{2}-2x+m\) realiosios nenulinės šaknys yra \(x_1,x_2\). Raskite visas \(m\) reikšmes, su kuriomis \mbox{\(\frac{x_1^2}{x_2^2}+\frac{x_2^2}{x_1^2}=7\)}.\end{minipage}
\item \begin{minipage}[t]{\linewidth}Duoto daugianario \(P(x)=x^{2}+2x-1\) šaknys yra \(x_1<x_2\). Nerasdami šaknų, apskaičiuokite \mbox{\(x_{1}^{2} - x_{2}^{2}\)}.\end{minipage}
\item \begin{minipage}[t]{\linewidth}Duoto daugianario \(P(x)=x^{2}-2x-1\) šaknys yra \(x_1,x_2\). Nerasdami jų, sudarykite kvadratinę lygtį, kurios vyriausiasis koeficientas lygus \(1\), o šaknys \mbox{\(\frac{x_{1}^{2}}{x_{2}},\quad \frac{x_{2}^{2}}{x_{1}}\)}.\end{minipage}
\item \begin{minipage}[t]{\linewidth}Raskite visas \(m\) reikšmes, su kuriomis daugianariai \[P(x)=2x^{2}+\left(m + 2\right)x+m - 2,\qquad Q(x)=x^{2}+x-1\] turi bent vieną bendrą realiąją šaknį.\end{minipage}
\item \begin{minipage}[t]{\linewidth}Raskite visas \(m\) reikšmes, su kuriomis lygties \(x^{2}-mx- m - 2=0\) realiosios šaknys \(x_1,x_2\) tenkina \mbox{\(x_1^3+x_2^3=0\)}.\end{minipage}
\item \begin{minipage}[t]{\linewidth}Lygties \(x^{2}-2x-1=0\) šaknys yra \(x_1,x_2\). Nerasdami šaknų, apskaičiuokite \mbox{\(x_{1}^{5} + x_{2}^{5}\)}.\end{minipage}
\item \begin{minipage}[t]{\linewidth}Kvadratinio daugianario \(P(x)\) vyriausiasis koeficientas lygus \(-1\), o realiosios šaknys tenkina: \[x_1^2+x_2^2=7,\quad x_1^3+x_2^3=-10.\] Raskite visus galimus daugianarius \(P(x)\).\end{minipage}
\item \begin{minipage}[t]{\linewidth}Lygties \(x^{2}+x-1=0\) šaknys yra \(x_1,x_2\). Nerasdami šaknų, apskaičiuokite \mbox{\(\frac{x_{1}^{2} + 1}{x_{1}^{3} + 2} + \frac{x_{2}^{2} + 1}{x_{2}^{3} + 2}\)}.\end{minipage}
\item \begin{minipage}[t]{\linewidth}Lygties \(x^{2}-\left(m + 1\right)x+1=0\) realiosios šaknys yra \(x_1,x_2\). Raskite visas \(m\) reikšmes, su kuriomis \mbox{\((x_1+1)^2,\quad(x_2+1)^2\)} yra lygties \(x^{2}-\left(4 m + 7\right)x+\left(m + 3\right)^{2}=0\) šaknys.\end{minipage}
\item \begin{minipage}[t]{\linewidth}Raskite visas \(m\) reikšmes, su kuriomis daugianario \(P(x)=x^{2}-\left(m + 1\right)x+3 m + 1\) abi šaknys yra sveikieji skaičiai.\end{minipage}
\item \begin{minipage}[t]{\linewidth}Daugianario \(P(x)=x^{2}-mx-1\) realiosios nenulinės šaknys yra \(x_1,x_2\). Raskite visas \(m\) reikšmes, su kuriomis \mbox{\(\left(\frac{x_1}{x_2}-\frac{x_2}{x_1}\right)^2=5\)}.\end{minipage}
\newpage
\item \begin{minipage}[t]{\linewidth}Raskite visas \(m\) reikšmes, su kuriomis daugianariai \[P(x)=x^{2}-mx+m + 1,\qquad Q(x)=x^{2}-3x+2\] turi bent vieną bendrą realiąją šaknį.\end{minipage}
\item \begin{minipage}[t]{\linewidth}Lygties \(2x^{2}-2x-1=0\) šaknys yra \(x_1<x_2\). Nerasdami šaknų, apskaičiuokite \mbox{\(x_{1}^{2} x_{2} - x_{1} x_{2}^{2}\)}.\end{minipage}
\item \begin{minipage}[t]{\linewidth}Duoto daugianario \(P(x)=x^{2}+x-1\) šaknys yra \(x_1,x_2\). Nerasdami šaknų, apskaičiuokite \mbox{\(x_{1}^{8} + x_{2}^{8}\)}.\end{minipage}
\item \begin{minipage}[t]{\linewidth}Raskite visas \(m\) reikšmes, su kuriomis lygties \(x^{2}-\left(m + 2\right)x+2 m=0\) realiosios šaknys \(x_1,x_2\) tenkina \mbox{\(x_1^3+x_2^3=35\)}.\end{minipage}
\item \begin{minipage}[t]{\linewidth}Lygties \(x^{2}+x-1=0\) šaknys yra \(x_1,x_2\). Nerasdami jų, sudarykite kvadratinę lygtį, kurios vyriausiasis koeficientas lygus \(1\), o šaknys \mbox{\(\left(x_{1} + 1\right)^{2},\quad \left(x_{2} + 1\right)^{2}\)}.\end{minipage}
\item \begin{minipage}[t]{\linewidth}Kvadratinio daugianario \(P(x)\) vyriausiasis koeficientas lygus \(1\), o realiosios šaknys tenkina: \[x_1^2+x_2^2=5,\quad x_1^4+x_2^4=17,\quad x_1+x_2>0.\] Raskite visus galimus daugianarius \(P(x)\).\end{minipage}
\item \begin{minipage}[t]{\linewidth}Lygties \(2x^{2}-mx+2=0\) realiosios šaknys yra \(x_1,x_2\). Raskite visas \(m\) reikšmes, su kuriomis \mbox{\(\frac1{x_1^2},\quad\frac1{x_2^2}\)} yra lygties \(x^{2}-\left(m - 1\right)x+1=0\) šaknys.\end{minipage}
\item \begin{minipage}[t]{\linewidth}Raskite visas \(m\) reikšmes, su kuriomis daugianario \(P(x)=2x^{2}-2 mx+4 m - 2\) abi šaknys yra sveikieji skaičiai.\end{minipage}
\item \begin{minipage}[t]{\linewidth}Duoto daugianario \(P(x)=2x^{2}-2x-1\) šaknys yra \(x_1,x_2\). Nerasdami jų, sudarykite kvadratinę lygtį, kurios vyriausiasis koeficientas lygus \(1\), o šaknys \mbox{\(\frac{1}{x_{1}^{2}},\quad \frac{1}{x_{2}^{2}}\)}.\end{minipage}
\item \begin{minipage}[t]{\linewidth}Raskite visas \(m\) reikšmes, su kuriomis daugianariai \[P(x)=x^{2}-\left(m + 1\right)x+2 m + 1,\qquad Q(x)=x^{2}-x-1\] turi bent vieną bendrą realiąją šaknį.\end{minipage}
\item \begin{minipage}[t]{\linewidth}Raskite visas \(m\) reikšmes, su kuriomis lygties \(-x^{2}+\left(m - 1\right)x+1=0\) realiosios šaknys \(x_1,x_2\) tenkina \mbox{\(x_1^3+x_2^3=4\)}.\end{minipage}
\item \begin{minipage}[t]{\linewidth}Duoto daugianario \(P(x)=-x^{2}+x+1\) šaknys yra \(x_1<x_2\). Nerasdami šaknų, apskaičiuokite \mbox{\(- \frac{1}{x_{2}} + \frac{1}{x_{1}}\)}.\end{minipage}
\item \begin{minipage}[t]{\linewidth}Kvadratinio daugianario \(P(x)\) vyriausiasis koeficientas lygus \(1\), o realiosios šaknys tenkina: \[(x_1-x_2)^2=5,\quad x_1^3+x_2^3=4.\] Raskite visus galimus daugianarius \(P(x)\).\end{minipage}
\item \begin{minipage}[t]{\linewidth}Lygties \(2x^{2}-2x-1=0\) šaknys yra \(x_1,x_2\). Nerasdami šaknų, apskaičiuokite \mbox{\(x_{1}^{6} + x_{2}^{6}\)}.\end{minipage}
\item \begin{minipage}[t]{\linewidth}Duoto daugianario \(P(x)=2x^{2}-2x-1\) šaknys yra \(x_1,x_2\). Nerasdami šaknų, apskaičiuokite \mbox{\(\frac{x_{1}^{2}}{x_{1}^{3} + 1} + \frac{x_{2}^{2}}{x_{2}^{3} + 1}\)}.\end{minipage}
\item \begin{minipage}[t]{\linewidth}Daugianario \(P(x)=2x^{2}-\left(m + 2\right)x+2\) realiosios nenulinės šaknys yra \(x_1,x_2\). Raskite visas \(m\) reikšmes, su kuriomis \mbox{\(\frac{x_1^2}{x_2}+\frac{x_2^2}{x_1}=18\)}.\end{minipage}
\end{enumerate}
\clearpage
\section*{Atsakymai}
\begingroup\small
\setlength{\baselineskip}{11pt}
\setlist[enumerate]{itemsep=0.02cm}
\begin{enumerate}
\item \(\frac{6}{7}\).
\item \(m=3\).
\item \(11\).
\item \(m=3 - \sqrt{3}\quad\text{arba}\quad m=\sqrt{3} + 3\).
\item \(m=\frac{5 - \sqrt{5}}{2}\quad\text{arba}\quad m=\frac{\sqrt{5} + 5}{2}\).
\item \(m=-6\quad\text{arba}\quad m=10\).
\item \(\begin{aligned}[t]P(x)=2x^{2}-4x-9,\\\text{arba}\quad P(x)=2x^{2}-10x+12.\end{aligned}\)
\item \(m=\frac{1 + \sqrt{17}}{2}\).
\item \(x^{2}-16x+46=0\).
\item \(- \frac{22 \sqrt{13}}{3}\).
\item \(m=-2\quad\text{arba}\quad m=3\).
\item \(x^{2}-4x-16=0\).
\item \(- 8 \sqrt{5}\).
\item \(P(x)=x^{2}-x-3\).
\item \(m=10\).
\item \(m=1\).
\item \(\frac{2}{3}\).
\item \(m=- \frac{\sqrt{2}}{2}\quad\text{arba}\quad m=\frac{\sqrt{2}}{2}\).
\item \(18\).
\item \(m=-4\quad\text{arba}\quad m=2\).
\item \(\frac{14}{17}\).
\item \(m=0\quad\text{arba}\quad m=2\).
\item \(m=-6\).
\item \(m=-4\quad\text{arba}\quad m=\frac{4}{5}\).
\item \(4 \sqrt{2}\).
\item \(x^{2}+14x-1=0\).
\item \(m=0\).
\item \(m=0\).
\item \(82\).
\item \(P(x)=-x^{2}-x+3\).
\item \(-1\).
\item \(m=2\).
\item \(m=-3\quad\text{arba}\quad m=13\).
\item \(m=-1\quad\text{arba}\quad m=1\).
\item \(m=5\).
\item \(\frac{\sqrt{3}}{2}\).
\item \(47\).
\item \(m=3\).
\item \(x^{2}-3x+1=0\).
\item \(\begin{aligned}[t]P(x)=x^{2}-x-2,\\\text{arba}\quad P(x)=x^{2}-3x+2.\end{aligned}\)
\item \(m=2 \left(1 + \sqrt{2}\right)\).
\item \(m=0\quad\text{arba}\quad m=8\).
\item \(x^{2}-8x+4=0\).
\item \(m=-3 - \sqrt{5}\quad\text{arba}\quad m=-3 + \sqrt{5}\).
\item \(m=2\).
\item \(- \sqrt{5}\).
\item \(P(x)=x^{2}-x-1\).
\item \(\frac{13}{2}\).
\item \(\frac{2}{3}\).
\item \(m=4\).
\end{enumerate}
\endgroup
\end{document}
